%%=============================================================================
%% CAPITULO 2 -- REFERENCIAL TEORICO
%%
%% Esta e a parte do trabalho que mais demanda dominio das regras de citacao
%% (diretas e indiretas). Reveja a NBR 10520 antes de escrever.
%%
%% Texto dummy (exercicio ia-6.1): trechos do pre-projeto do autor (versao 7).
%%=============================================================================

\section{Referencial teórico}

O referencial teórico será organizado em torno de uma questão central: a geração
de dados sintéticos não deve ser avaliada apenas pela capacidade de produzir
registros plausíveis, mas pela relação entre fidelidade, utilidade analítica e
risco de exposição. Essa perspectiva permite deslocar o foco de uma avaliação
centrada exclusivamente no algoritmo gerador para uma avaliação orientada ao uso
do artefato, coerente com o caráter aplicado da pesquisa.

\subsection{Dados sintéticos e o desafio dos dados tabulares}

Para dados tabulares heterogêneos, a escolha do modelo torna-se particularmente
importante. \citeonline{xu2019} destacam que tabelas combinam variáveis discretas
e contínuas, distribuições multimodais e categorias desbalanceadas,
características que dificultam a aplicação direta de modelos tradicionais. A
CTGAN foi proposta como uma arquitetura condicional para enfrentar parte desses
problemas, utilizando condicionamento e estratégias específicas para variáveis
categóricas.

Uma das contribuições mais importantes para a avaliação de dados sintéticos é a
distinção entre utilidade geral e utilidade específica proposta por
\citeonline{snoke2018}. A utilidade geral procura avaliar se os dados sintéticos
possuem distribuição comparável à dos dados originais; a utilidade específica,
por sua vez, compara os resultados de uma análise concreta realizada nos dois
datasets. Essa distinção é fundamental porque um dataset pode reproduzir
distribuições marginais de forma satisfatória e, ainda assim, falhar em preservar
relações importantes para uma tarefa analítica.

\subsection{Avaliação da utilidade de dados sintéticos}

A necessidade de relacionar a métrica à finalidade da avaliação aparece no
estudo de validação de métricas de utilidade conduzido por El Emam e
colaboradores:

\begin{citacao}
A regular task by developers and users of synthetic data generation (SDG)
methods is to evaluate and compare the utility of these methods. Multiple
utility metrics have been proposed and used to evaluate synthetic data.
However, they have not been validated in general or for comparing SDG methods.
\parencite{elemam2022}
\end{citacao}

No estudo, foram avaliadas seis métricas em 30 datasets e três métodos de
geração, e a distância multivariada de Hellinger baseada em uma representação
Gaussian Copula apresentou o melhor desempenho para o objetivo de ranquear os
métodos \parencite{elemam2022}. O resultado reforça uma questão central desta
pesquisa: a métrica escolhida deve ser relacionada à finalidade da avaliação, e
não apenas à facilidade de cálculo.
